\section{Discussion}\label{sec:discussion}
\begin{table}[t]\centering\footnotesize
\caption{Pre-registered hypotheses of the first two rounds (H1--H5, H6--H10), the estimate with its 95\% interval and
the verdict. (a) marks amendments made before the results; (p) marks readings or tests chosen after the results.
Contrasts are candidate minus comparator. Contrasts across batches are raw; seed-adjusted values are in Table~\ref{tab:batches}.}\label{tab:hyp}
\setlength{\tabcolsep}{3pt}
\begin{tabular}{>{\raggedright\arraybackslash}p{0.30\linewidth}>{\raggedright\arraybackslash}p{0.47\linewidth}>{\raggedright\arraybackslash}p{0.16\linewidth}}\toprule
hypothesis (as registered) & estimate & verdict \\ \midrule
H1 placebo loops show a positive selection gain of predictable size & first-round best scores inside leave-one-out
predictive intervals; six loops $\vhPlaceboMin$--$\vhPlaceboMax\%$ vs $\vhPlaceboPred\%$ predicted & supported \\
H2 DGM loops' selection gain exceeds their fresh gain; fresh gain positive & per-loop excess (tasks resampled)
$\vhtwoMinInfl$ to $\vhtwoMaxInfl$ points, every interval includes 0; six loops (hierarchical) $\vhtwoMean$
[$\vhtwoLo$, $\vhtwoHi$]; fresh gains all positive & not supported \\
H3 best log-free edit reaches $\geq 0.8\times$ the loops' fresh gain & (a) on $F_1$: ratio $\vhThreeRatio$ & not
supported \\
H4 from the engineered seed, fresh gains of all arms within $\pm 3$ points, selection gains $> 0$ & DGM loops:
\nhFourCI{} on $F$ (p: 90\% two one-sided tests) and $+\ngainStrongB$ [$\ngainStrongBlo$, $\ngainStrongBhi$] on $F_1$;
single-edit arms ($F_1$): $\nladSCone$ [$\nladSConelo$, $\nladSConehi$] and $\nladSCtwo$ [$\nladSCtwolo$,
$\nladSCtwohi$]; selection gains of the loops $+\nSgainStrongA$ and $+\nSgainStrongB$ & not supported (equivalence not
shown) \\
H5 evolved-agent gains shrink or vanish under Haiku~4.5 and Sonnet~4.6 & (a) on $F_1$, loops 1 and 3 (loop 2's agent
not run); change of gain, loop 3's agent: $\vhfiveKThree$ [$\vhfiveKThreelo$, $\vhfiveKThreehi$], $\vhfiveSThree$
[$\vhfiveSThreelo$, $\vhfiveSThreehi$]; loop 1's: $\vhfiveKOne$ [$\vhfiveKOnelo$, $\vhfiveKOnehi$], $\vhfiveSOne$
[$\vhfiveSOnelo$, $\vhfiveSOnehi$] & supported only for loop 3 under Sonnet~4.6 \\ \midrule
H6 best textbook combination $\geq$ best first-round loop agent $-3$ (and cheaper) & (a) chosen on $F_1$
(\vHsixWinner), on $F_2$: $\vHsixFtwo$ [$\vHsixFtwolo$, $\vHsixFtwohi$]; literal, best on all of $F$ (\vhsixLitBest; chosen and scored on the same problems):
$\vhsixLit$ [$\vhsixLitlo$, $\vhsixLithi$], seed-adjusted $\vhsixLitAdj$ [$\vhsixLitAdjlo$, $\vhsixLitAdjhi$]; (a)
pre-specified single candidates on all of $F$: all four moves $\vpsEBLPRaw$ [$\vpsEBLPRawlo$, $\vpsEBLPRawhi$], L
$\vpsLRaw$ [$\vpsLRawlo$, $\vpsLRawhi$]; all cheaper & not supported; literal reading supported on raw scores, not
after seed adjustment \\
H7 random-diagnosis loops $\geq$ first-round DGM loops $-3$ (mean fresh gain) & $\vhRandomMinusDgm$ [$\vhRandomMinusDgmlo$,
$\vhRandomMinusDgmhi$]; (p) vs six DGM loops $\vhRandomMinusDgmAll$ [$\vhRandomMinusDgmAlllo$, $\vhRandomMinusDgmAllhi$]
& not supported \\
H8 best of five one-shot redesigns $\geq$ best loop agent $-3$; mean redesign $\geq$ mean loop best $-3$ & (a) best on
$F_2$: $\vHeightOpusbest$ [$\vHeightOpusbestlo$, $\vHeightOpusbesthi$]; literal, best on all of $F$:
$\vheightLitOpusRaw$ [$\vheightLitOpusRawlo$, $\vheightLitOpusRawhi$]; mean: $\vHeightOpusmean$ [$\vHeightOpusmeanlo$,
$\vHeightOpusmeanhi$] & best: not supported; mean: supported \\
H9 two edits carry $\geq 80\%$ of a loop's fresh gain & loop 1 has exactly two edits; for loop 2 half and for loop 3
all but one of the leave-one-out variants do not apply & not testable (loops 2, 3) \\
H10 seed $+$ textbook moves ($+$ retry) $\geq$ published best $-3$ on fresh tasks & (a; pilot results known) DGM
SWE-bench with Sonnet~4.5, moves: $\vsweMovesMinusBest$ [$\vsweMovesMinusBestlo$, $\vsweMovesMinusBesthi$] as run,
$\vsensMovesMinusBest$ [$\vsensMovesMinusBestlo$, $\vsensMovesMinusBesthi$] with early-ended patches kept; moves and
retries: $\vsweRetryMinusBest$ [$\vsweRetryMinusBestlo$, $\vsweRetryMinusBesthi$] and $\vsensRetryMinusBest$
[$\vsensRetryMinusBestlo$, $\vsensRetryMinusBesthi$]; ADAS (CoT-SC vs agent): holds on MGSM, DROP, MMLU, fails on
GPQA; HGM and DGM Polyglot not run & (a) supported for DGM SWE-bench; 3 of 4 for ADAS \\ \bottomrule
\end{tabular}\end{table}

\begin{table}[t]\centering\footnotesize
\caption{Pre-registered hypotheses of the third round (H11--H15), the estimate with its 95\% interval and the
verdict. Every contrast is within one batch.}\label{tab:hyp3}
\setlength{\tabcolsep}{3pt}
\begin{tabular}{>{\raggedright\arraybackslash}p{0.30\linewidth}>{\raggedright\arraybackslash}p{0.47\linewidth}>{\raggedright\arraybackslash}p{0.16\linewidth}}\toprule
hypothesis (as registered) & estimate & verdict \\ \midrule
H11 at equal output budget, seed $+$ textbook moves $\geq$ DGM best agent $-3$ (SWE-bench) & as run
$\wMovesBest$ [$\wMovesBestlo$, $\wMovesBesthi$]; early-ended patches kept $\wsMovesBest$ [$\wsMovesBestlo$,
$\wsMovesBesthi$] & supported (also superior) \\
H12 at equal output budget, DGM best agent vs its initial agent (expected: no significant gain) & as run $\wBestInit$
[$\wBestInitlo$, $\wBestInithi$]; patches kept $\wsBestInit$ [$\wsBestInitlo$, $\wsBestInithi$] & no gain; worse as
run \\
H13 (descriptive) textbook agent vs initial agent with the raised limit & $\wMovesInit$ [$\wMovesInitlo$,
$\wMovesInithi$] & -- \\
H14 clean random-diagnosis loops $\geq$ all DGM loops $-3$ (mean fresh gain) & $\rcHfourteen$ [$\rcHfourteenlo$,
$\rcHfourteenhi$] & \rcHfourteenVerdict \\
H15 log-free redesigns: $\geq 4$ of 5 raise the cap; mean gain $\geq$ log-based $-3$ & \lfRaised{} of 5;
$\lfDiff$ [$\lfDifflo$, $\lfDiffhi$] & supported \\ \bottomrule
\end{tabular}\end{table}

Tables~\ref{tab:hyp} and~\ref{tab:hyp3} list the registered hypotheses and their outcomes, and the findings in the introduction
summarize them together with the exploratory comparisons. Read together, they show that the loops mostly
found standard engineering, which a one-shot redesign also reaches at a lower cost per task whether or not it sees
failure logs; they do not show that loops find nothing.

Two readings follow. The first concerns what search contributes. In most of our settings, a baseline that uses no
search (textbook moves, or a single redesign from the same failure logs) reached the level of the loops' best agents,
or came within a few points of it, at lower cost. The exceptions run both ways: on gpt-oss-120b an agent found by a loop
on another model gained more than our textbook set did, in a comparison across batches, and ADAS's discovered agent wins on GPQA.
Diagnosis also carries information: once the request list contains nothing aimed at the empty-turn failure,
random-diagnosis loops gain much less than DGM loops (\S\ref{sec:random-clean}). The loops' costs are
many evaluations and inflated selection scores (\S\ref{sec:round1}). The second reading concerns transfer to a
new task model. A harness is tuned to the model it was built for, whether by a loop or directly: the engineered seed
loses to the DGM seed under gpt-oss-120b, and it is ahead of loop 3's agent under Haiku~4.5 and Sonnet~4.6, about
level under Sonnet~5.5, and behind it under gpt-oss-120b and Gemini~2.5 Flash (Appendix~\ref{app:round1}). Published
self-improvement results should therefore be re-measured when the task model changes.

\paragraph{Recommendations.} Claims of self-improvement should report (i) gains on tasks that no selection step saw,
(ii) a placebo loop or its order-statistics prediction, (iii) the seed resampled at the agent's cost, (iv) the seed
with textbook moves and a one-shot redesign from the same failure logs, measured in the same batch as the loop's agent
(our own registered tests did not fully meet this), and (v) the released agent re-run with a newer model. Harnesses
should keep hidden tests out of the agent's container, keep the agent's work when its process ends early, and check that
each task's reference patch still passes.

\paragraph{Limitations.} Loops were run on one task family with one to nine loops per arm, and the confidence intervals
over problems hold the realized loops fixed. At the level of loops, differences between loop types are significant only
for the clean random-diagnosis loops of the third round (\S\ref{sec:random-clean}). The SWE-bench and other-model evaluations use one run per task, and budget and time limits censor several
agents under other models (Appendix~\ref{app:cens}). Claude~3.5 Sonnet, the DGM's original model, is no longer served,
so we cannot separate the effect of the model change from a failure to replicate. We could not determine what ends the
best agent's runs early under Claude Sonnet~4.5; machine load and memory do not explain it (\S\ref{sec:swe-equal}). The SWE-bench batch at equal output budget used \wN{} tasks. HGM's released agent and the DGM's Polyglot agents were not run. The proxy caps output at 16{,}000 tokens for
every arm. The engineered seed, the source of the textbook moves, was written while pilot runs that included \vnpN{}
of the fresh problems took place, and their logs showed the DGM seed's empty responses, so the textbook moves are not
independent of what failure logs show; without those problems the main numbers barely change (Appendix~\ref{app:pilot}).
The request list of the second-round random-diagnosis loops was drafted after the first round had exposed the
empty-turn failure, and three of its requests aim at or near it (\S\ref{sec:round2}); the third round repeats the test
without them (\S\ref{sec:random-clean}). Several
comparisons were added after earlier results (Appendix~\ref{app:dev}, items 3 and 15); those of item 3 were specified
before they were computed, and those of item 15 after they were computed. Under Haiku~4.5 most runs of the DGM seed end with its agent process killed by a signal, which the engineered
seed avoids (Appendix~\ref{app:othermodels}). Heavy load on the shared machine affected part of one exploratory batch
and the last generations of three loops, and the second-round SWE-bench run used an oversubscribed machine
(Appendix~\ref{app:dev}).
